%=============================================================================!
% 18.0  CHAPTER OPENING                                                       !
%=============================================================================!
A spring balance in a shop is calibrated by hanging known weights on it
and marking where the pointer stops. Each reading carries a small error,
from the pointer's width and the eye that reads it, so no single weight
fixes the scale; the marks are drawn so that they miss the readings as
little as possible, taken together. A learned interatomic potential is
made the same way. Its readings are energies and forces from a
first-principles calculation, its marks are the numbers inside the model,
and drawing them is a search for the numbers that miss the readings
least. Chapter~17 ran such a model, MACE-MP-0, as a black box; this
chapter begins to open the box; the planned chapters that follow develop
the network models themselves.

We begin with a loss, a number that measures how far a model misses its
readings. A least-squares fit first determines one weight, the stiffness
of a bond, then several weights in a curve through noisy readings of
argon's pair energy. This exposes a difficulty: enough adjustable
weights can fit the noise as well as the underlying curve. A penalty on
large weights controls that freedom, and its interpretation as a prior
leads to Gaussian processes and their uncertainty estimates.

When a fit cannot be obtained in one calculation, we lower the loss by
gradient descent. Its step has a stability limit like an integrator's;
mini-batches, momentum and Adam change the cost and progress of that
search. The final sections separate the readings used to train a model
from those used to judge it, showing why consecutive trajectory frames
must not be split at random. Planned Chapter~19 builds neural networks
and trains them with these optimisers.
